\documentclass[11pt,a4paper]{article}
\usepackage[utf8]{inputenc}
\usepackage[T1]{fontenc}
\usepackage[spanish]{babel}
\usepackage{hyperref}
\usepackage{geometry}
\geometry{margin=2.5cm}
\usepackage{parskip}

\title{004 el miedo}
\author{Alejandro Valenzuela Rivera\\\small Universidad Autónoma de Nuevo León}
\date{\today}
\begin{document}
\maketitle

\begin{quote}
\textbf{Sesgo declarado.} Este trabajo se sostiene desde el sesgo declarado de su autor, y no finge neutralidad. Lo que afirma se puede tumbar, y ese es el punto: el falsificador está escrito, no escondido.
\end{quote}

\section*{Resumen}
El miedo no se pregunta si se quita: se corre contra la cuenta; sin datos es un resultado.

\section{Por qué se sostiene}
El miedo es un proceso de costo y no una virtud.

\section{De dónde viene}
Sinestro, Green Lantern.

\section{El caso}
La rata de las cavernas: sobreestimar sobrevive, subestimar muere.

\section{Método}
Una decision que cuesta algo, y que se toma con el costo declarado.

\section{Límites, y qué lo tumba}
Una decision sin costo: no habia nada en juego.

\section{Lo que queda abierto}
\textquestiondown{}Que le pasa a una especie sin miedo?

\section*{Cómo se leyó esta tesis}
Esto se escribió en el repositorio man-made-empathy, del que este candado es una pieza: candados/02 el alma y el libre albedrio/004 el miedo.md. La máquina que lo compuso lo hizo desde este mismo archivo, y no agregó nada que no estuviera escrito. Si algo de esto es falso, se dice en el mismo sitio, y se corrige con fecha.

\end{document}
